\chapter{Literature Review and Market Research}
\label{ch:literature}

\begin{guidance}[What this chapter does]
Establishes what has already been done, what problems remain and which
technologies could solve yours. The review must \emph{justify} your aims and
objectives and lead into the specification. Typical length: 1,500--2,500
words.

A review is not a list of summaries. Group sources by theme, compare them,
say which are reliable and why, and finish every section with what it means
for \emph{your} project.

Useful sources (Handbook, Phase Two): academic journals and conference
papers, standards (British Standards Online via the Library), patents,
manufacturers' data, market reports and statistical data. Cite everything in
Harvard Solent style with \verb|\textcite{key}| or \verb|\parencite{key}|.
\end{guidance}

\section{Search method}
\label{sec:search}

\begin{guidance}
Briefly state where and how you searched: databases (e.g.\ IEEE Xplore,
ScienceDirect, Google Scholar, the Solent Library), keywords, and how you
selected sources. Two or three sentences make the review reproducible.
\end{guidance}

\section{Theme 1: \ldots}
\label{sec:theme1}

Several approaches to modelling the thermal behaviour of buildings
represent walls and air volumes by equivalent electrical networks
\parencite{achterbosch1985}. \todo{Your first theme.}

\section{Theme 2: \ldots}
\label{sec:theme2}

\section{Existing products and market}
\label{sec:market}

\begin{guidance}
Competitive and analogous products, with their key figures in a comparison
table. State where the data came from (datasheet, website with the date
viewed).
\end{guidance}

\begin{table}[htbp]
  \centering
  \caption{Comparison of existing products (example)}
  \label{tab:market}
  \begin{tabular}{lrrl}
    \toprule
    Product   & Price (£) & Accuracy (\si{\celsius}) & Interface \\
    \midrule
    Product A & 45  & $\pm 0.5$ & USB       \\
    Product B & 120 & $\pm 0.1$ & Wi-Fi     \\
    Product C & 30  & $\pm 1.0$ & SD card   \\
    \bottomrule
  \end{tabular}
\end{table}

\section{Summary and implications for the project}
\label{sec:lit-summary}

\begin{guidance}
The most important section of the chapter: state the gap or the opportunity
your project addresses, and how the findings informed the specification
in \cref{ch:specification}.
\end{guidance}
